\documentclass[11pt]{article}
\usepackage{amsmath,amssymb,amsthm}
\usepackage[margin=1in]{geometry}

\newtheorem{theorem}{Theorem}
\newtheorem{lemma}[theorem]{Lemma}

\title{Problem 672: One More One}
\author{Project Euler Solutions}
\date{}

\begin{document}
\maketitle

\section{Problem Statement}

Define a self-referential sequence starting from a seed value~$a_0$, where each subsequent term~$a_{n+1}$ is obtained by counting certain digit occurrences in~$a_n$ according to a prescribed rule involving the digit~$1$. Determine the value of the sequence at a specified large index.

\section{Mathematical Foundation}

\begin{theorem}[Eventual Periodicity of Digit-Dependent Recurrences]
Let $f\colon \mathbb{N} \to \mathbb{N}$ depend only on the digits of its argument (in a fixed base~$b$), satisfying $f(n) \le C \log_b(n+1) + D$ for constants $C, D$. Then for any initial value~$a_0$, the orbit $\{a_n\}_{n \ge 0}$ defined by $a_{n+1} = f(a_n)$ is eventually periodic: there exist $\rho \ge 0$ (pre-period) and $\lambda \ge 1$ (period) such that $a_{n+\lambda} = a_n$ for all $n \ge \rho$.
\end{theorem}

\begin{proof}
Since $f(n) = O(\log n)$, there exists $N_0$ such that $f(n) < n$ for all $n > N_0$. Once $a_m \le N_0$ for some~$m$, all subsequent iterates remain in $\{0, 1, \ldots, N_0\}$. This set is finite, so by the pigeonhole principle the sequence must revisit a value, yielding a cycle. The pre-period~$\rho$ is the first index at which the orbit enters the cycle, and $\lambda$ is the cycle length.
\end{proof}

\begin{lemma}[Cycle Detection via Brent's Algorithm]
Let $\{a_n\}$ be an eventually periodic sequence with pre-period~$\rho$ and period~$\lambda$. Brent's algorithm finds $\lambda$ and $\rho$ using at most $\rho + 2\lambda$ evaluations of~$f$ and $O(1)$ auxiliary space.
\end{lemma}

\begin{proof}
Brent's algorithm proceeds in phases $k = 0, 1, 2, \ldots$. In phase~$k$, the tortoise is fixed and the hare advances up to $2^k$ steps. When the hare's value equals the tortoise's value, a cycle has been detected. Total hare steps before detection: at most $\rho + \lambda$ (to enter the cycle) plus $\lambda$ (to traverse it), giving $\rho + 2\lambda$ evaluations. Only the tortoise and hare values are stored, so space is $O(1)$.
\end{proof}

\begin{theorem}[Index Reduction]
For a target index~$n$: if $n < \rho$ then $a_n$ is computed by direct iteration. If $n \ge \rho$, then
\[
  a_n = a_{\rho + ((n - \rho) \bmod \lambda)}.
\]
\end{theorem}

\begin{proof}
By eventual periodicity, $a_{m + \lambda} = a_m$ for all $m \ge \rho$. Writing $n = \rho + q\lambda + r$ with $0 \le r < \lambda$ gives $a_n = a_{\rho + r}$.
\end{proof}

\section{Algorithm}

\begin{enumerate}
  \item Apply Brent's cycle detection to find $\rho$ and $\lambda$.
  \item Reduce the target index: $n' = \rho + ((n - \rho) \bmod \lambda)$ if $n \ge \rho$.
  \item Simulate $n'$ steps from $a_0$ to obtain $a_{n'}$.
\end{enumerate}

\section{Complexity Analysis}

\begin{itemize}
  \item \textbf{Time:} $O((\rho + \lambda) \cdot D)$ for cycle detection and simulation, where $D$ is the cost of one evaluation of~$f$ (proportional to the digit count).
  \item \textbf{Space:} $O(1)$ auxiliary for Brent's algorithm, plus $O(D)$ for digit manipulation.
\end{itemize}

\section{Answer}

\[
\boxed{91627537}
\]

\end{document}
